\section{Introduction and motivation}
\label{sec:intro}

The MINERvA medium-energy (ME) FHC charged-current (CC) inclusive
double-differential cross section in muon transverse and longitudinal momentum,
$d^{2}\sigma/(d\pT\,d\pPar)$, was published in arXiv:2106.16210 (Ruterbories
\emph{et al.}) using D'Agostini iterative Bayesian unfolding (IBU) on a binned
response matrix. This note reproduces that measurement with \emph{unbinned}
OmniFold, an unfolding technique based on iterated reweighting with machine-learned
likelihood ratios that operates on per-event observables rather than on a fixed
histogram binning.

\paragraph{Why unbinned OmniFold.}
\begin{itemize}
  \item \textbf{No binning at unfold time.} OmniFold learns event weights in the
        full observable space; the result is histogrammed only at the end, so the
        binning can be chosen (or refined) after unfolding without re-running.
  \item \textbf{Natural path to higher dimensions.} Adding observables is adding
        input features, not multiplying response-matrix bins --- this is the
        enabling property for the available-energy extension
        (\S\ref{sec:openquestions}) that has no practical IBU analogue.
  \item \textbf{Different regularisation.} OmniFold's regulariser is the
        iteration count and the model capacity of the gradient-boosted decision
        trees (GBDTs), not the IBU prior-iteration coupling. Comparing the two
        on the same data quantifies the method-dependence of the result
        (\S\ref{sec:results}).
\end{itemize}

\paragraph{Goal and scope.}
The goal is a faithful, fully uncertainty-quantified reproduction of the
published 2D cross section as the demonstrator for unbinned OmniFold on MINERvA
data, followed by the available-energy higher-dimensional extension. This note
covers the method (\S\ref{sec:method}), the 2D result and its comparison to the
published data and to GENIE MINERvA Tune v1 (\S\ref{sec:results}), the
uncertainty budget (\S\ref{sec:systematics}), the validation suite
(\S\ref{sec:validation}), and the open questions for the collaboration
(\S\ref{sec:openquestions}).
